% ------------------------------------------------------------------------------
% §7 Mechanism — drafted Phase C from:
%   interpretation_staged_framework.md (+ amendments); findings_job2_contamination.md;
%   findings_task1.md / task1b.md (post-fix); findings_task5b.md; findings_task9/10;
%   Exhibits: T3
% ------------------------------------------------------------------------------
\section{Mechanism: A Standing Posture, Not a Daily Calibration}\label{sec:mechanism}

The regression results say that withholding deepens with the compensation
price where the scheme's eligibility bites. This section opens the conduct up
at the level of individual bids, days, and direction episodes. The picture
that emerges is specific: the absence is a \emph{standing posture} --- chosen
at the day-block level, indifferent to the daily security envelope and to the
daily prize --- and the compensation mechanism pays gross rates into that
posture whenever the system's need makes the unit essential. The section ends
with the two findings that discipline any stronger claim: a state-dependence
result this project itself withdrew after a contamination audit, and the
grain at which the dose response does and does not appear.

\subsection{What a direction buys, and what it pays}

Across the 740 direction episodes with a resolved in-force bid, directed
output equals the unit's operating floor almost exactly: median directed
output 40~MW per Torrens unit --- the floor block itself --- with 55.9\% of
episodes delivering zero or negative excess over the floor and 80.8\% within
25~MW (Figure~\ref{fig:floor}). A direction buys presence, not energy.

It pays as if it bought energy at scarcity prices. Compensation reconciles to
$0.95 \times \text{directed MWh} \times \dt$ with $R^{2}=0.99$
(Section~\ref{sec:setting}), and the aggregate consequence over 2022--2024 is
the mechanism's headline economics: the direction channel paid the three
Torrens units \$141.9M for 625,635 directed MWh --- their 683 episodes of
the 740 above --- that would have earned \$7.2M
at spot prices over the same windows --- a \$134.7M gap, a ratio of roughly
20 to 1 that holds for every unit in every year. In 2023 the market value of the directed output
was \emph{negative} for all three units: 48.2\% of directed megawatt-hours
were delivered into negative spot prices. Nor is the comparison distorted by
averaging: at the interval level, spot exceeds marginal fuel cost in 3.9\% of
directed intervals, the median directed interval loses \$141/MWh at fuel cost
alone, and only 11 of 683 episodes (1.6\%) were profitable on average. There
is no commercial running inside these windows to forgo. Given these relative
payment rates, a standing absent posture is not a puzzle to be explained ---
it is the only stance an informed desk would hold.

\subsection{The posture is standing, not state-dependent}

Three tests --- each with its interpretation committed in the project
record before estimation --- all on bids formed outside direction exposure
(the clean-day rule of Section~\ref{sec:design}), locate the choice; the
broader $N-1$ essentiality flag they require is constructed and validated in
Appendix~\ref{app:n1}.

\emph{The base rate.} The units' day-ahead bids withhold or price out the
floor megawatts on 76--80\% of all clean days --- essential or not. Absence is
the resting state, held in multi-day blocks (median unavailability spell 2--4
days, maxima up to 100 days), total rather than partial (when the units withdraw
they declare zero, not a derating: sub-floor positive declarations are 0.2\%
of intervals), and organised as a rotating one-to-two-of-three station roster
(all three Torrens units are on offer together on 4\% of days, and unit
availabilities are deliberately anti-correlated).

\emph{The exit act is event-timed, not need-timed.} On clean days when a
Torrens unit has an evening on offer, cancelling it is \emph{less} likely ---
by 16 percentage points off a 22.5\% base ($p=0.02$, bootstrap $p=0.017$;
Table~\ref{tab:mechanism}, column 1) --- when the system is one contingency
from needing the unit. A loss-timing account fails alongside: cancellations
are less likely, not more, when the previous day's prices imply running
losses (column 2, wrong sign at $p=0.005$). The evening exits that precede
directions sit in the direction run-up itself, not in the security calendar
or the loss calendar.

\emph{The price of presence does not move either.} When the units offer their
floor megawatts at all, they offer them at the floor-band price ---
essentially $-\$1{,}000$ --- in essential hours, in ordinary hours, and on
ordinary days alike: the essential-day effect on the offered floor price is
+\$45 ($p=0.86$; column 4), and Pelican Point's within-day essential-hour
effect is exactly zero. The pricing margin
joins the availability margin: the posture is set by the day-block and does
not respond to the envelope.

\begin{table}[!htbp]\centering
\caption{Mechanism: the exit act is state-dependent; floor pricing is not}
\label{tab:mechanism}
\begin{adjustbox}{max width=\textwidth}
\begin{threeparttable}
\begin{tabular}{lcccc}
\toprule
 & \multicolumn{3}{c}{Exit act (evening offered share)} & Floor price (\$/MWh) \\
\cmidrule(lr){2-4}\cmidrule(lr){5-5}
 & (1) & (2) & (3) & (4) \\
\midrule
Essential ($N{-}1$) & $-0.162^{**}$ & $-0.124^{*}$ & & $45.3$ \\
 & (0.066) & (0.065) & & (265.3) \\
Essential ($N{-}1$) $\times$ expected loss & & $-0.001^{***}$ & & \\
 & & (0.000) & & \\
$N{-}1$ only (not $N{-}0$) & & & $-0.190^{***}$ & \\
 & & & (0.061) & \\
Directly pivotal ($N{-}0$, \textit{pex}) & & & $-0.001$ & \\
 & & & (0.106) & \\
\midrule
Observations & 473 & 473 & 473 & 650 \\
\bottomrule
\end{tabular}
\begin{tablenotes}\footnotesize
\item Day-level regressions on clean days (bids formed outside direction exposure), $N{-}1$ essentiality flag (Task~7 construction). Columns (1)--(3): outcome is the evening offered share (exit act); column (2) adds the expected foregone-loss control and its interaction; column (3) splits essentiality into tiers. Column (4): outcome is the offered floor price. Wild-cluster-bootstrap-based inference in the source analysis confirms column (1) at $p=0.017$. Standard errors clustered by month in parentheses. Both tests' interpretations were committed in the project record before estimation; the $N{-}1$ flag construction is Appendix~\ref{app:n1}.
\item $^{*}p<0.10$, $^{**}p<0.05$, $^{***}p<0.01$ (analytic).
\end{tablenotes}
\end{threeparttable}
\end{adjustbox}
\end{table}


\begin{figure}[!htbp]\centering
\includegraphics[width=0.8\textwidth]{fig4_output_vs_floor}
\caption{Directed output relative to the unit's operating floor, for the 740
direction episodes with a resolved in-force bid, 2022--2024. The mass at zero
is the institutional design: a direction compels presence at minimum stable
load.}
\label{fig:floor}
\end{figure}

\subsection{A withdrawn result: the contamination audit}

An earlier day-grain version of the essentiality test in this project found
what looked like clean state-dependence: units withdrew 12.7 percentage
points more on essential days ($p<0.001$). The finding did not survive its
own audit, and the correction is reported here as part of the result. On 54\%
of essential days, the day-ahead bid was formed while a direction was running
or pending; a bid formed under direction mechanically shows availability. On
the 80 essential days with cleanly formed bids the withdrawal effect is
$-0.015$ ($p=0.71$): the apparent state-dependence was the direction itself
sitting in the measured bids. What survives on the clean record is the
standing posture --- and the fact that 95.4\% of the 280 clean-day first
directions arrived at a day-ahead stance that had already withdrawn (71.8\%)
or priced out (23.6\%) the floor megawatts. The system directs into a posture
that predates its need.

\subsection{Reconciling the grains}

The clean-day, day-grain dose-response test is null (49 essential days over
14 months; all interactions $p=0.50$--$0.96$), while the interval-grain
matched test of Section~\ref{sec:results} finds $-0.051$ per \$100/MWh at
$p<0.01$. These are not in tension; they are different instruments pointed at
the same object. The day-grain test conditions on clean essential days ---
which the contamination audit shows are scarce precisely because the
direction regime blankets essential periods --- and its wide confidence
intervals cannot exclude effects several times the interval estimate. The
interval-grain test pools 12,513 essential unit-intervals against matched
comparisons with 300 times the observations. The honest summary is a powered
estimate and an uninformative bound, not a conflict. What the day-grain
record does add is the shape of the response: the posture does not flex
day-to-day with the envelope or the prize (this section), while its depth
across months tracks what essentiality pays (Section~\ref{sec:results}).
Both are what a policy chosen once --- and paid gross whenever the system's
need arrives --- would produce.\footnote{The daily activity around the
standing choice is stage-appropriate housekeeping, not repositioning:
document-revision activity ramps into directions (a vigilance signature), but
net availability and price-ladder movement around direction issue is
approximately zero. A desk that watches the weather and a desk that watches
the direction calendar are watching the same clock in this system; the staged
pattern shows the desk is informed, and this paper deliberately does not use
it to establish motive. The economic content is in the mechanism's design, not in generator
sophistication.}
